
\begin{table}[H]
\centering
\caption{Thực nghiệm bổ sung trên DB15K:
lựa chọn teacher và distillation; MRR (\%).}
\label{tab:dsom_ablation}
\scriptsize
\setlength{\tabcolsep}{3pt}
\renewcommand{\arraystretch}{1.12}

\begin{tabular}{@{}lll>{\columncolor{resultgray}}r@{}}
\toprule
\textbf{Run}
& \textbf{Teacher selection}
& \textbf{KD}
& \textbf{MRR$\uparrow$} \\
\midrule

Exp01 & Frozen RC & NDKD
& \dsomlink{gvpp95st}{41.30} \\

A1 & Structure & NDKD
& \dsomlink{3fhw62i2}{41.27} \\

A2 & Description & NDKD
& \dsomlink{jxfadxq5}{38.18} \\

A3 & S + D & NDKD
& \dsomlink{sf7vexg5}{41.37} \\

A4 & Average & NDKD
& \dsomlink{e6i6aw58}{41.31} \\

A5 & Learned RC & NDKD
& \dsomlink{qqnzvnz2}{41.42} \\

A6 & RC (continuous CE reward) & NDKD
& \dsomlink{zuhqedof}{41.29} \\

A7 & S + D & Vanilla KD
& \dsomlink{i9azqqx2}{\textbf{41.44}} \\

A8 & S + D & DKD
& \dsomlink{opus98xu}{41.40} \\

A9 & Learned RC & Vanilla KD
& \dsomlink{9vw2fikt}{41.34} \\

A10 & Learned RC & DKD
& \dsomlink{hpruwc92}{41.30} \\

\bottomrule
\end{tabular}

\par\smallskip

\begin{minipage}{0.98\linewidth}
\scriptsize
Exp01 giữ nguyên policy RC, trong khi A5
huấn luyện policy bằng optimizer riêng.
A3/A7/A8 sử dụng tổ hợp teacher cố định S+D;
A5/A9/A10 sử dụng learned RC.
\end{minipage}


\end{table}


A5 và A6 đều dùng NDKD, cùng checkpoint teacher và các đặc
trưng đầu vào; student được tối ưu bằng Adagrad với learning
rate $0.1$, còn policy dùng Adam riêng với learning rate
$10^{-5}$. Khác biệt cấu hình chính là
\texttt{rc\_reward\_mode}: A5 dùng \texttt{binary}, còn A6 dùng
\texttt{continuous\_ce}.
\footnote{\href{https://github.com/keiranoluv/Reproduce_DSoM/blob/ab1ee9b25f5b3385b5fe0e01d1cfed1f497db18e/scripts/run_db15k_a5_learned_rc_full.sh\#L69-L77}{Cấu hình A5};
\href{https://github.com/keiranoluv/Reproduce_DSoM/blob/ab1ee9b25f5b3385b5fe0e01d1cfed1f497db18e/scripts/run_db15k_a6_continuous_ce_full.sh\#L69-L77}{cấu hình A6}.}
A5 dùng reward $+1$ nếu CE của tổ hợp teacher thấp hơn
student và $-10$ trong trường hợp còn lại. A6 dùng reward
liên tục và baseline:
\begin{align*}
 R_i &= \mathrm{CE}_{\mathrm{stu},i}
        -\mathrm{CE}_{\mathrm{sel},i},\\
 \widehat{R}_i &= \mathrm{CE}_{\mathrm{stu},i}
        -\mathrm{CE}_{\mathrm{avg},i}.
\end{align*}
Vì vậy, advantage của A6 là
$\delta_i=\mathrm{CE}_{\mathrm{avg},i}
-\mathrm{CE}_{\mathrm{sel},i}$, với ``avg'' là trung bình logits
của ba teacher.
\footnote{\href{https://github.com/keiranoluv/Reproduce_DSoM/blob/ab1ee9b25f5b3385b5fe0e01d1cfed1f497db18e/src/kd.py\#L500-L532}{Triển khai reward và advantage của A5/A6.}}
Đây là biến thể reward do nhóm bổ sung. Trên một seed,
MRR của A5 và A6 lần lượt là $41.42\%$ và $41.29\%$;
kết quả này chưa đủ để kết luận về ưu thế ổn định giữa hai
cách tính reward.
